\uuid{JEO6}
\titre{Changements de variables}
\auteur{hanani ; mzari}
\organisation{exo7}
\licence{CC BY-NC-SA 4.0}
\datecreate{2020-01-20}
\niveau{L1}
\module{Analyse}
\chapitre{Calcul d'intégrales}
\sousChapitre{Autre}
\difficulte{2}
\theme{Calculs d'intégrales, changements de variables}
\exo7qcmid{lille-2.integrales.010}

\contenu{
\begin{question}
\qclassification{127.99}
\qtags{motcle=changements de variables}
Parmi les affirmations suivantes, cocher celles qui sont vraies :
\begin{answers}
\bad{$\displaystyle \int _{0}^{\pi /2}\cos ^2x\sin x\,\mathrm{d}x=-\frac{1}{3}$.}
\good{$\displaystyle \int _{0}^{\pi /2}\sin ^2x\cos x\,\mathrm{d}x=\frac{1}{3}$.}
\bad{$\displaystyle \int _1^4\frac{\mathrm{e}^{\sqrt{x}}}{\sqrt{x}}\mathrm{d}x=\frac{\mathrm{e}^2}{2}-\mathrm{e}$.}
\good{$\displaystyle \int _{-\pi/2}^{\pi/2}\frac{\cos x\,\mathrm{d}x}{2+\sin x}=\ln 3$.}
\end{answers}
\vskip2mm
\begin{explanations}
Avec $u=\cos x\Rightarrow \mathrm{d}u=-\sin x\,\mathrm{d}x$ : $\displaystyle \int _0^{\pi/2}\cos ^2x\sin x\,\mathrm{d}x=-\int _1^0u^2\, \mathrm{d}u=\frac{1}{3}$.
\vskip0mm
Avec $u=\sin x$, on a : $\mathrm{d}u=\cos x\,\mathrm{d}x$ et
$\displaystyle \int _0^{\pi/2}\sin ^2x\cos x\,\mathrm{d}x=\int _0^1u^2\, \mathrm{d}u=\frac{1}{3}$.
\vskip0mm
Ensuite, avec $u=\sqrt{x}$, on a : $\displaystyle \mathrm{d}u=\frac{\mathrm{d}x}{2\sqrt{x}}$ et
$\displaystyle \int _1^4\frac{\mathrm{e}^{\sqrt{x}}}{\sqrt{x}}\mathrm{d}x=2\int _1^2\mathrm{e}^{u}\mathrm{d}u=2(\mathrm{e}^2-\mathrm{e})$.
\vskip0mm
Enfin, avec $u=2+\sin x$, on obtient : $\mathrm{d}u=\cos x\, \mathrm{d}x$ et $\displaystyle \int _{-\pi/2}^{\pi/2}\frac{\cos x\,\mathrm{d}x}{2+\sin x}=\int _{1}^{3}\frac{\mathrm{d}u}{u}=\ln 3$.
\end{explanations}
\end{question}
}
